% ===================================================================
%  ТАБЛИЦЯ № 11 — Місто і його будівлі. — Пожежа.
%  Traduction 2026 d'apres texto/io, controlee sur texto/fr, texto/en
%  et texto/es. Consigne : voir texto/uk/10-tablycia-01.tex.
%
%  OUVERTURE DE LA TROISIEME SERIE : « ТРЕТЯ СЕРІЯ ». « третя » etait
%  entre dans ORDINALO avec les trois autres ordinaux ukrainiens ; le
%  role serio se pose tout seul.
%
%  LE BLOC c1-03-1 EST CELUI QUI A COUTE UNE INVERSION AU FINNOIS.
%  « la pordo (8) dil olda fortifikita kastelo (7) » : le possede porte
%  le numero, et une langue qui prepose son genitif renverse les deux
%  renvois. L'ukrainien postpose le sien, comme l'ido —
%  « брама (8) старого укріпленого замку (7) » — et ne paie rien.
% ===================================================================

%%K t11-apar-1 apar
\VUsaut{2.0mm}
\VUcentre{13.2pt}{90}{ТРЕТЯ СЕРІЯ}
\VUsaut{3.0mm}
\VUfilet{16mm}
\VUsaut{5.6mm}
\VUtitre{15.0pt}{60}{ТАБЛИЦЯ № 11}
\VUsaut{1.4mm}
\VUpk{11.4pt}{40}{Місто і його будівлі. --- Пожежа.}
\VUsaut{2.2mm}

%%K t11-tit-1 sub
\VUpk{10.2pt}{40}{Околиці.}
\VUsaut{3.0mm}

%%K t11-c1-01-1 p
1. --- Я живу у великому місті за сто кілометрів від океану, яке проте є
першорядним \VUgras{портом} \textsuperscript{(1)}. Річка, що тече вздовж нього,
завширшки чотириста метрів і утворює прекрасний рейд.

%%K t11-c1-02-1 p
2. --- Уся частина міста вздовж \VUgras{набережних} \textsuperscript{(2)} має
вигляд цілком морський і промисловий і утворює \VUgras{передмістя}
\textsuperscript{(3)}, де живуть самі моряки та робітники. Останні працюють на
\VUgras{заводах} \textsuperscript{(4)} і на \VUgras{газовому заводі}
\textsuperscript{(5)}, чий високий \VUgras{димар} \textsuperscript{(6)} і
\VUgras{газгольдер} видно здалеку.

%%K t11-c1-03-1 p
3. --- У Середні віки місто було укріплене, і деякі рештки його валів ще
знаходять, надто \VUgras{браму} \textsuperscript{(8)} старого \VUgras{укріпленого
замку} \textsuperscript{(7)} з двома \VUgras{баштами} \textsuperscript{(9)} по
боках, які тепер правлять за \VUgras{в'язницю}.

%%K t11-c1-04-1 p
4. --- У всьому цьому \VUgras{кварталі}, на вигляд дуже старому, тільки одна
будівля порівняно нова: \VUgras{митниця} \textsuperscript{(10)}. Вона стоїть між
набережною і маленькою \VUgras{вулицею} \textsuperscript{(11)}, що веде до старої
\VUgras{ратуші} \textsuperscript{(12)}. Цю останню будівлю легко впізнати по
велетенській \VUgras{дозорній вежі} \textsuperscript{(13)}, яка височіє над
старим містом.

%%K t11-c1-05-1 p
5. --- По другий бік вулиці \VUgras{стоїть} будівля в тому самому стилі, що й
митниця: це \VUgras{біржа} \textsuperscript{(14)}. Один з її фасадів виходить на
маленьку площу, де міститься \VUgras{префектура} \textsuperscript{(15)}. Наше
місто, правда, не столиця, але все ж важливий повітовий центр.

%%K t11-tit-2 sub
\VUsaut{5.0mm}
\VUpk{10.2pt}{40}{Горішня частина міста.}
\VUsaut{3.4mm}

%%K t11-c1-06-1 p
6. --- Залога, доволі велика, міститься в просторих і дуже здорових
\VUgras{касарнях} \textsuperscript{(16)}; вони тягнуться до самих ґрат
\VUgras{міського парку} \textsuperscript{(17)}. \VUgras{Бульвар}
\textsuperscript{(19)}, що веде до цього парку, проходить позаду \VUgras{ощадної
каси} \textsuperscript{(18)} і виводить на площу \VUgras{собору}
\textsuperscript{(20)}.

%%K t11-c1-07-1 p
7. --- Усі ці будівлі підіймаються уступами на невеликому пагорбі, де стоїть і
наш великий \VUgras{ліцей} \textsuperscript{(21)}.

%%K t11-c1-07-2 p
Якщо спуститися злегка похилою вулицею, що виходить на Велику вулицю, дійдеш до
\VUgras{суду} \textsuperscript{(22)}, перед яким лежить приємний \VUgras{сквер}
\textsuperscript{(26)}. \VUgras{Площа} ця невелика, але посеред міста вона
утворює оазу зелені й прохолоди. Тому її охоче відвідують городяни, які люблять
посидіти під повіткою \VUgras{кав'ярні} \textsuperscript{(25)}, слухаючи
військовий оркестр, що грає щочетверга і щонеділі в \VUgras{альтанці} \textsuperscript{(27)}
серед газонів і клумб.

%%K t11-c1-08-1 p
8. --- На одному з цих газонів стоїть бронзова \VUgras{статуя}
\textsuperscript{(28)}, пам'ятник, поставлений на спомин про одного з наших
городян, який зробив великі послуги рідному місту.

%%K t11-tit-3 sub
\VUsaut{4.0mm}
\VUpk{10.2pt}{40}{Пересування.}
\VUsaut{3.4mm}

%%K t11-c1-09-1 p
9. --- Наше місто ще не освітлене електрикою; у нього тільки \VUgras{газові
ліхтарі} \textsuperscript{(29)}. Але \VUgras{місцева преса} клопочеться про те,
щоб цей спосіб освітлення було поліпшено, \VUgras{так само як} уже поліпшено
\VUgras{тягу} на трамваях. \VUgras{Старі} вагони, \VUgras{які тягли}
\VUgras{коні}, замінено зручними \VUgras{електричними трамваями}
\textsuperscript{(31)}, які швидко возять пасажирів з одного кінця \VUgras{міста
в другий} за дуже \VUgras{низьку} плату.

%%K t11-c1-10-1 p
10. --- Біля суду --- \VUgras{банк} \textsuperscript{(32)}, де враховують
торговельні векселі. \VUgras{Банкові посильні} \textsuperscript{(33)} ходять по
домах збирати належне.

%%K t11-tit-4 sub
\VUsaut{4.0mm}
\VUpk{10.2pt}{40}{Мистецтво і наука.}
\VUsaut{3.4mm}

%%K t11-c1-11-1 p
11. --- Гарна будівля неподалік --- \VUgras{музей}. У ньому містяться разом
міська \VUgras{бібліотека} \textsuperscript{(34)}, прекрасні \VUgras{природничі
збірки} \textsuperscript{(35)} (природничо-історичний музей) і зграбна
\VUgras{картинна галерея} \textsuperscript{(36)}, що складає чудову постійну
\VUgras{виставку}.

%%K t11-c1-11-2 p
Він відкритий для публіки двічі на тиждень, але відвідувачі можуть оглянути
його будь-якого дня за невеликий дарунок \VUgras{сторожеві}
\textsuperscript{(37)}. \VUgras{Художники} \textsuperscript{(38)} приходять туди
працювати. Їх бачиш \VUgras{перед} мольбертами, з \VUgras{палітрою} в руці, як
вони копіюють славетні картини.

%%K t11-c1-12-1 p
12. --- На висотах за музеєм ледве розрізняєш \VUgras{баню}
\textsuperscript{(40)} \VUgras{лікарні} \textsuperscript{(39)} і будівлі
\VUgras{учительської семінарії} \textsuperscript{(41)}, де вчилися викладачі
наших міських шкіл. На Театральній площі є \VUgras{пошта}
\textsuperscript{(43)}, вміщена в \VUgras{приватному домі}
\textsuperscript{(42)}.

%%K t11-tit-5 sub
\VUsaut{5.0mm}
\VUpk{11.4pt}{40}{Пожежа.}
\VUsaut{3.0mm}

%%K t11-tit-6 sub
\VUpk{10.2pt}{40}{Крик «Пожежа!».}
\VUsaut{3.4mm}

%%K t11-c2-01-1 p
1. --- Містом розноситься тривожна вість: у театрі щойно почалася пожежа! Коли я
добираюся до \VUgras{площі} \textsuperscript{(96)}, натовп уже зібрався на
\VUgras{тротуарі} \textsuperscript{(46)} і на \VUgras{бруківці}
\textsuperscript{(47)}. \VUgras{Поштові урядовці} \textsuperscript{(44)} і
\VUgras{листоноші} \textsuperscript{(45)} виходять із пошти.
\VUgras{Вагоновод} \textsuperscript{(48)} змушений був спинити свій трамвай, щоб
пропустити міську \VUgras{карету швидкої допомоги} \textsuperscript{(50)}, яка
прибула з \VUgras{хірургом} \textsuperscript{(52)} і \VUgras{доглядальницею}
\textsuperscript{(51)}, щоб допомогти \VUgras{пораненому} \textsuperscript{(53)},
якого на \VUgras{ношах} \textsuperscript{(54)} несуть до місця біля
\VUgras{газетного кіоска} \textsuperscript{(55)}.

%%K t11-c2-02-1 p
2. --- Рота піхоти, що верталася з навчання, спинилася, щоб помогти підтримувати
лад разом із кінними \VUgras{міськими вартовими} \textsuperscript{(61)}.
\VUgras{Вершники}, чиї коні стають дибки, дають раду краще за \VUgras{солдатів}
\textsuperscript{(57)} і \VUgras{жандармів} \textsuperscript{(63)}, відтісняючи
\VUgras{роззяв} \textsuperscript{(56)}. Жандарми до того ж дуже зайняті. Один із
них указує \VUgras{поліціянтам} \textsuperscript{(64)}, куди нести іншого
\VUgras{потерпілого} \textsuperscript{(65)}, хороброго пожежника, який упав з
драбини.

%%K t11-c2-03-1 p
3. --- \VUgras{Офіцер} \textsuperscript{(66)} указує достеменно, куди поставити
\VUgras{парову помпу} \textsuperscript{(67)}, яка якраз під'їжджає. Чекаючи на цю
потужну машину, він велів поставити \VUgras{ручну помпу} \textsuperscript{(68)},
довгий \VUgras{рукав} \textsuperscript{(69)} якої ллє на вогонь потоки води. Її
качають шестеро пожежників, а їхній товариш, тримаючи \VUgras{брандспойт}
\textsuperscript{(70)}, спрямовує \VUgras{струмінь} \textsuperscript{(71)} у саме
серце пожежі.

%%K t11-c2-04-1 p
4. --- По другий бік театру підняли велику \VUgras{драбину}
\textsuperscript{(72)}. Вона дозволяє пожежникам підібратися до полум'я, що
здіймається серед клубів чорного \VUgras{диму} \textsuperscript{(75)}.

%%K t11-tit-7 sub
\VUsaut{5.0mm}
\VUpk{10.2pt}{40}{Сміливі вчинки.}
\VUsaut{3.4mm}

%%K t11-c2-05-1 p
5. --- \VUgras{Вогонь} \textsuperscript{(74)} почався у фойє \VUgras{театру}
\textsuperscript{(73)}, під час репетиції \VUgras{опери}, перша \VUgras{вистава}
якої мала відбутися завтра. На щастя, в залі було лише кілька \VUgras{глядачів}
\textsuperscript{(76)}. Бідні люди сховалися на \VUgras{балконі}
\textsuperscript{(77)}, куди хоробрі \VUgras{пожежники} \textsuperscript{(86)}
йдуть їм на поміч із драбиною і мотузками.

%%K t11-c2-06-1 p
6. --- Під \VUgras{портиком} \textsuperscript{(78)}, де \VUgras{афіша}
\textsuperscript{(79)} оголошує виставу, видно, як біжать \VUgras{музиканти}
\textsuperscript{(81)} \VUgras{оркестру} \textsuperscript{(80)}, несучи свої
інструменти: \VUgras{скрипку} \textsuperscript{(81)}, \VUgras{флейту}
\textsuperscript{(82)}, \VUgras{віолончель} \textsuperscript{(83)},
\VUgras{тромбон} \textsuperscript{(84)}, \VUgras{барабан}
\textsuperscript{(85)} і так далі. Намагаються врятувати частину \VUgras{обстави}
\textsuperscript{(87)}.

%%K t11-c2-07-1 p
7. --- \VUgras{Акторам} \textsuperscript{(89)} і \VUgras{акторкам}
\textsuperscript{(88)}, \VUgras{співакам} і співачкам, довелося вибігти в
сценічних \VUgras{костюмах} \textsuperscript{(90)}. Тенор пояснює
\VUgras{журналістові} \textsuperscript{(92)}, як це сталося. \VUgras{Репортер}
прибіг із найпершої хвилини разом із владою: \VUgras{префектом}
\textsuperscript{(91)}, \VUgras{мером} \textsuperscript{(93)},
\VUgras{поліцмейстером} \textsuperscript{(94)} і \VUgras{слідчим}
\textsuperscript{(95)}, який поведе дізнання, щоб установити, на кому лежить
відповідальність.
\VUsaut{6.0mm}
\VUfilet{22mm}
